\documentclass[12pt]{article}
\usepackage{amsmath,amssymb,amsthm}
\usepackage{booktabs}
\usepackage{hyperref}
\usepackage[margin=1in]{geometry}

\newcommand{\E}{\mathbb{E}}
\newcommand{\Wtopo}{W_{\mathrm{topo}}}
\newcommand{\Winfo}{W_{\mathrm{info}}}

\title{The Blackwell Dilemma: Computational Appendix}
\author{Replication Package}
\date{}

\begin{document}
\maketitle

This document accompanies ``The Blackwell Dilemma: When Better Information Makes Agents Worse Off Under Endogenous Feasibility Constraints.'' We report results from six families of computational experiments, all conducted on 2D torus graphs with i.i.d.\ $\mathrm{Uniform}[0,1]$ rewards. Full simulation code and data are available alongside this document.

\section{Phase Transition Verification}

On $n = 50$ ($2{,}500$ vertices) with $M = 400$ Monte Carlo trials per parameter value: mean welfare loss increases from 0.32 at $p = 0$ to 0.47 at $p = 0.95$. The reachable fraction drops sharply at $p \approx 0.5$ (from $>0.99$ to $<0.08$), confirming the Harris--Kesten threshold $p_c = 1/2$. At $\beta = 3.0$, the bimodality coefficient (Sarle's $b$) crosses the $5/9$ threshold near $p_c$ and peaks above $b = 0.60$ in the supercritical regime, confirming the two-population prediction.

\section{Decomposition Verification}

The decomposition $W = \Wtopo + \Winfo$ is verified numerically ($n = 40$, $M = 400$). Five values of $\beta \in \{0.5, 1, 2, 5, 20\}$ are overlaid: all five $\Wtopo$ curves are visually identical, confirming $\partial \Wtopo / \partial \beta = 0$. The informational component $\Winfo$ is smooth, decreasing with $\beta$, and approximately zero at $\beta = 20$. The phase transition in total welfare comes entirely from the topological component.

\section{Counterfactual Experiment}

As a within-model counterfactual, we fix $\beta = 5.0$ and start at $p = 0.70$ (supercritical), then surgically repair blocked edges to lower effective $p$:
\begin{center}
\begin{tabular}{ccc}
\toprule
Effective $p$ & Mean Loss & Reachable Fraction \\
\midrule
0.70 (original) & 0.370 & 0.002 \\
$\approx 0.29$ (partial repair) & 0.192 & 0.99 \\
$\approx 0.00$ (full repair) & 0.146 & 1.00 \\
\bottomrule
\end{tabular}
\end{center}
Same cognition ($\beta = 5.0$), different topology $\to$ different outcome. The welfare floor disappears when $p$ crosses below $p_c$, confirming the loss is topological ($p$), not informational ($\beta$).

\section{Policy Gradient Comparison}

We compare the marginal effectiveness of structural policy ($\partial W / \partial p$) against information policy ($\partial W / \partial \beta$) via central finite differences ($\Delta p = 0.03$, $\Delta \beta = 0.3$; $n = 40$, $M = 400$ trials per configuration; bootstrap standard errors from $B = 2{,}000$ resamples):
\begin{center}
\begin{tabular}{ccccc}
\toprule
$(p, \beta)$ & $|\partial W/\partial p|$ & $|\partial W/\partial \beta|$ & Ratio & Interpretation \\
\midrule
$(0.45, 3.0)$ & $0.457 \pm 0.218$ & $0.011 \pm 0.022$ & $42\times$ & Near $p_c$: structural dominates \\
$(0.55, 5.0)$ & $0.683 \pm 0.294$ & $0.026 \pm 0.029$ & $26\times$ & Supercritical \\
$(0.70, 5.0)$ & $0.542 \pm 0.331$ & $0.038 \pm 0.031$ & $14\times$ & Deep supercritical \\
\bottomrule
\end{tabular}
\end{center}
At all three parameter configurations, $|\partial W/\partial \beta|$ is not significantly different from zero (point estimate $<$ bootstrap SE), while $|\partial W/\partial p|$ is significantly positive ($>2$ SE from zero at $p = 0.55$). Near $p_c$, structural policy is $42\times$ more effective than information policy. The ratio remains large ($>10\times$) throughout the supercritical regime.

\section{Finite-Size Scaling}

Grid sizes $n \in \{20, 40, 60, 80, 100\}$ (up to $10{,}000$ vertices) with $M = 800$ trials per parameter value. The transition sharpens monotonically with system size: at $p_c = 0.5$, the reachable fraction converges toward 0.5 from above (0.608 at $n=20$ to 0.491 at $n=100$), consistent with the critical point of 2D bond percolation. Data collapse using $\nu = 4/3$ yields good alignment for both the welfare loss and reachable fraction curves. The welfare loss at $p = 0.25$ is near zero for all grid sizes ($< 0.004$) while at $p = 0.75$ it stabilizes near 0.29, confirming size-independent supercritical loss. The variance of the welfare loss at $p_c$ is approximately size-independent ($\approx 0.04$), reflecting that the loss is a bounded $[0,1]$ functional of the percolation cluster rather than the unbounded susceptibility observable.

\section{Bimodality Emergence}

At $\beta = 3.0$, the welfare loss distribution transitions from unimodal (heavy-tailed, right-skewed) at $p < 0.3$ to visibly bimodal near $p_c$: a ``free'' mode at $\ell \approx 0$ and a ``trapped'' mode at $\ell \approx 0.25$. At lower precision ($\beta = 1.0$), the distribution remains unimodal at all $p$---both populations exist but are blurred by signal noise. This confirms that $\beta$ acts as a resolution parameter separating the two populations, as predicted by the theory.

\section{Analytic Formula Verification}

The theory predicts $\E[|\Wtopo| \mid |R| = k] = (n-k)/((n+1)(k+1))$ for i.i.d.\ Uniform$[0,1]$ rewards. We verify this formula via Monte Carlo: for each grid size $n \in \{20, 40\}$ and blocking probability $p \in \{0.3, 0.5, 0.6, 0.7\}$, we run $M = 2{,}000$ trials, bin outcomes by realized cluster size $|R| = k$, and compare the empirical conditional mean loss against the analytic prediction. In the supercritical regime ($p = 0.7$), where cluster sizes are small and the formula predicts $\Theta(1)$ losses, the median relative error across $k$-bins is $10.2\%$ ($n=20$) and $8.7\%$ ($n=40$), confirming the analytic formula. At $p = 0.6$, median errors are $21.6\%$ ($n=20$) and $16.3\%$ ($n=40$). In the subcritical regime ($p = 0.3$), both the analytic and empirical values are near zero ($<10^{-3}$) and relative errors are dominated by Monte Carlo noise, as expected.

\end{document}
